\chapter{文献综述}

在随机波动率模型的发展历程中，3/2 模型凭借其独特的非线性扩散结构与对极端市场波动的优异捕捉能力，逐渐成为金融工程领域期权定价与风险管理方向的研究热点。本章将遵循“模型范式演进 - 理论框架发展 - 定价方法迭代 - 市场应用拓展 - 国内研究现状”的逻辑主线，系统梳理 3/2 模型相关领域的国内外研究成果，最终通过文献述评明确现有研究的核心缺口，确立本研究的学术定位与边际贡献。

\section{随机波动率模型的演进}

期权定价理论的演进过程，本质上是市场参与者对金融资产收益率“非正态性”与波动率“时变风险”认知不断深化的过程。自 Black-Scholes-Merton（BSM）模型\citet{blackPricingOptionsCorporate1973,mertonTheoryRationalOption1973}确立了无套利定价的基本范式以来，常数波动率与标的价格服从几何布朗运动的核心假设，为衍生品定价提供了简洁且高效的解析框架，奠定了现代金融工程的基石。然而，随着全球金融市场复杂性的持续提升，BSM 模型的核心假设与真实市场特征产生了系统性背离。大量实证研究表明，全球各类资产的收益率分布普遍存在显著的“尖峰厚尾”特征，期权市场隐含波动率层面则呈现出随行权价格变化的“波动率微笑”与随到期期限变化的“波动率偏斜”现象\citet{rubinsteinNonparametricTestsAlternative1985}，这一现象在市场危机与极端行情中表现得尤为突出，揭示了市场对尾部风险溢价的非线性定价特征，也推动了波动率模型从确定性向随机性的范式转换。

为修正 BSM 模型的理论缺陷，学术界首先发展出局部波动率模型。\citet{dermanRidingSmile1994}与\citet{Dupire1994}分别从离散图与连续偏微分方程视角，构建了依赖于标的价格与时间的确定性局部波动率框架，能够完美拟合某一时刻的隐含波动率曲面。但局部波动率模型本质上仍属于确定性波动率范畴，无法刻画波动率自身的随机波动特征，对未来波动率曲面的预测能力较弱，难以适配动态对冲与长期衍生品定价的实践需求。

在此背景下，随机波动率（Stochastic Volatility, SV）模型应运而生，其核心思想是将波动率本身设定为一个随机扩散过程，从而同时刻画标的价格的随机波动与波动率自身的时变特征。以\citet{hestonSimpleNewFormula1997}为代表的仿射随机波动率模型（Affine SV Models），通过引入平方根扩散过程（也称 1/2 模型）刻画波动率的均值回归属性，推导标的对数价格特征函数的仿射闭式解，极大地降低了模型校准与期权定价的计算复杂度，成为全球衍生品行业的定价标准基石。随后，针对利率与外汇市场的复杂波动动态，\citet{Hagan2002}提出的 SABR 模型通过引入随机波动率项与波动率 - 标的价格的相关结构，进一步强化了对隐含波动率曲线形态的捕捉能力，成为场外利率衍生品定价的主流模型。

然而，随着市场微观结构的演变与极端行情的频繁出现，传统 1/2 模型的内生局限性逐渐凸显。学者们发现，Heston 模型在拟合极超短期期限权的隐含波动率陡峭时存在理论不足，其线性漂移项难以解释波动率在极端危机时刻的爆发式增长与非线性均值回归特征\citep{jonesDynamicsStochasticVolatility2003}。这驱动了学术界从仿射框架向非仿射框架的探索。3/2 模型作为非仿射随机波动率模型的典型代表，其扩散项与瞬时方差的 3/2 次方成正比，漂移项呈现出非线性的均值回归结构，这种更强的非线性特征赋予了模型更灵活的峰度刻画能力与尾部风险捕捉能力，使其在处理资产价格剧烈跳变与波动率高度聚集的场景时，展现出比传统 1/2 模型更深厚的理论底蕴与更优的实证表现。

\section{3/2 模型及其扩展框架的发展}

自 \citet{lewisOptionValuationStochastic2000} 对 3/2 随机波动率模型的基础解析性质进行系统性推导以来，学术研究围绕 3/2 模型的理论框架优化、解析性质拓展与数值实现方案开展了大量研究，形成了较为完备的理论成果体系。

在模型基础解析框架的构建方面，\citet{carrNewApproachOption2007} 开发了新颖的定价框架，通过直接对方差互换率进行建模，绕过了对瞬时波动率风险价格建模的复杂过程，为 3/2 模型的风险中性定价提供了全新的实现路径。针对 3/2 模型作为非仿射模型带来的期权定价与校准难题，\citet{medvedevApproximationCalibrationShortTerm2007} 提出了基于小期限渐近展开的综合分析框架，推导出非仿射局部波动率模型下隐含波动率的闭式展开式，为包括 3/2 模型在内的非仿射模型的快速截面校准提供了核心理论工具，也成为后续 3/2 模型实证研究的重要校准方法基础。

在模型结构的扩展与优化方面，为实现股票现货与 VIX 衍生品的一致性建模，\citet{baldeauxConsistentModelingVIX2012}在纯扩散 3/2 模型的基础上引入同步跳跃项，构建了“带跳跃的 3/2 模型”（3/2 plus jumps model），进一步提升了模型对极端市场极端波动的刻画能力。\citet{grasselli42Stochastic2017} 在其开创性研究中提出了“4/2 随机波动率模型”，将瞬时方差的扩散项设定为 1/2 项与 3/2 项的线性组合，从而将 Heston 模型与 3/2 模型有机地整合在同一框架内，形成了通用的随机波动率模型框架，有效结合了两类模型的优势，极大地拓展了 3/2 模型的理论应用边界。针对 3/2 模型的非仿射特性带来的模拟难题，\citet{kouarfateExplicitSolutionSimulation2021} 基于 3/2 模型方差过程的显式解式，开发了一套全新的模拟框架，该框架巧妙利用了 3/2 模型的方差过程是 CIR 过程倒数的核心数学性质，为 3/2 模型的蒙特卡洛模拟提供了新的实现思路。

\section{精确模拟方法的演进}

在随机波动率模型的定价实践中，蒙特卡洛模拟是处理路径依赖型衍生品、高维模型定价的核心工具，而模拟方法的核心发展方向，是在消除时间离散化误差的同时，兼顾计算效率与全场景数值稳定性。自精确模拟方法被提出以来，学术界围绕仿射模型与非仿射 3/2 模型的精确模拟方案开展了持续的优化研究。

在仿射随机波动率模型的精确模拟奠基性研究中，\citet{broadieExactSimulationStochastic2006} 首次提出了 Heston 随机波动率模型及其他仿射跳跃扩散过程下，标的股票价格与方差过程的无偏精确模拟方法，被称为精确蒙特卡洛（Exact MC）格式。该方法通过对积分方差的条件分布进行抽样，彻底消除了时间离散化带来的系统误差，为随机波动率模型的精确模拟奠定了理论基础。随后，\citet{glassermanGammaExpansionHeston2011} 对 Broadie-Kaya 算法进行了关键优化，通过利用 \citet{pitmanDecompositionBesselBridges1982} 关于扩散过程积分级数的分解结论，将条件积分方差分解为无限个伽马随机变量的级数形式，避免了原算法中计算成本极高的数值拉普拉斯逆变换，显著提升了精确模拟的计算效率。

针对非仿射 3/2 模型的精确模拟，\citet{baldeauxEXACTSIMULATION322012} 开展了开创性研究，首次探讨了 3/2 模型下股票价格过程的精确模拟方案。该研究利用 \citet{craddockCalculationExpectationsClasses2009} 基于李对称分析得出的扩散过程转移密度结论，成功将 \citet{broadieExactSimulationStochastic2006} 针对仿射过程的精确模拟算法改编并应用于 3/2 模型，推导出 3/2 模型下积分方差的条件拉普拉斯变换闭式解，为 3/2 模型的无偏模拟提供了首个理论可行的实现方案。

近年来，3/2 模型的精确模拟算法在数值稳定性与计算效率上取得了进一步突破。\citet{choiSimulationSchemesHeston2024} 提出了一种无需使用高维修正贝塞尔函数的新型精确模拟方案，其核心发现是当用于模拟终端方差的泊松变量为条件时，3/2 模型的条件积分方差可以得到显著简化，彻底解决了原 Baldeaux 算法中修正贝塞尔函数在极端参数下易出现数值溢出的核心痛点。此外，\citet{brignoneExactSimulationStochastic2026} 提出了一种利用条件傅里叶余弦（COS）方法的新型精确模拟框架，通过余弦级数展开构造积分方差的条件累积分布函数，实现了积分方差的高效高精度抽样，解决了自 \citet{broadieExactSimulationStochastic2006} 以来精确模拟中长期存在的误差控制难、计算开销大的实施难题，进一步拓展了精确模拟方法的工程应用边界。

\section{傅里叶变换方法的研究}

傅里叶变换定价方法是随机波动率模型下欧式期权快速定价与全面校准的核心技术路径，其核心思想是将期权价格的风险中性期望转化为标的对数价格特征函数的复数域积分，通过数值积分算法实现快速求解，完全规避了蒙特卡洛模拟的统计误差与时间离散化误差。

在傅里叶变换定价方法的奠基性研究中，\citet{carrOptionValuationUsing1999} 提出了基于快速傅里叶变换（Fast Fourier Transform, FFT）的期权定价方法，通过引入衰减因子解决了期权支付函数的傅里叶变换不可积问题，推导出欧式期权价格的傅里叶逆变换闭式表达式，结合 FFT 算法实现了在 $O(N \log N)$ 时间复杂度内一次性计算全行权价区间的期权价格，成为随机波动率模型下期权定价的主流工程方法。针对 FFT 方法存在的栅栏效应、仅能计算等距行权价格的局限，\citet{fangNovelPricingMethod2009} 提出了基于傅里叶余弦级数展开的 COS 定价方法，通过将标的收益率的概率密度函数在有效支撑区间上展开为余弦级数，推导出期权价格的闭式级数求和形式，实现了指数级的收敛速度，同时可灵活计算任意行权价对应的期权价格，在非仿射随机波动率模型中展现出极强的适用性与计算效率优势。

对于 3/2 模型而言，由于其非仿射的模型结构，标的对数价格并不具备仿射模型中简单的指数仿射特征函数，其特征函数的解析表达涉及复杂的特殊函数，为傅里叶变换方法的应用带来了一定挑战。\citet{lewisOptionValuationStochastic2000} 关于复数域积分定价的早期研究，为 3/2 模型的解析性质探索提供了核心理论框架；\citet{dufresne2001integrated} 的研究进一步指出，3/2 过程的积分分布与 Whittaker 函数及修正贝塞尔函数密切相关，完整推导了 3/2 模型下积分方差的分布性质，为特征函数的闭式求解奠定了数学基础。在实务应用中，FFT 定价法被广泛应用于 3/2 模型的参数校准与期权定价，但其在极短到期时间、极端波动率参数下仍面临数值稳定性不足、计算耗时过长的挑战\citet{drimusOptionsRealizedVariance2011}。而 COS 方法凭借其指数级数收敛速度与强数值稳定性，在 3/2 模型的定价与校准中展现出显著的性能优势，成为近年来该领域的研究重点。

\section{加密货币期权市场中随机波动率模型的研究}

随着以比特币（BTC）和以太坊（ETH）为代表的加密资产进入机构化交易时代，加密货币期权市场规模呈现指数级增长，其独特的市场动力学特征也成为金融工程领域的前沿研究方向。不同于传统权益、外汇市场，加密货币市场具有 $7\times24$ 小时不间断交易、无涨跌幅限制、内生波动率极高、极端行情频发的核心特征，其收益率分布呈现出显著的尖峰厚尾特征与极高的波动率的波动率（Vol-of-Vol, VoV），对传统期权定价模型提出了严峻挑战。

\citet{houPricingCryptocurrencyOptions2020}的研究指出，加密资产的波动率展现出比传统股票资产更强的持久性和更高的波动率的波动率，市场极端尾部风险显著高于传统金融市场，这使得 BSM 模型甚至基础的 Heston 模型在加密期权定价中经常出现严重的估值偏差，尤其是在极端行情下，模型对波动率微笑的拟合能力大幅下降。针对这一现象，学术界主要从两个方向探索适配加密市场的模型框架：一类研究通过在随机波动率模型中引入跳跃扩散过程（Jump-Diffusion），捕捉加密货币价格的瞬时极端跳变\citet{scaillet2020high}；另一类研究则侧重于改进波动率过程的非线性结构，探索非仿射随机波动率模型在加密市场的适配性。

近年来，由于 3/2 模型在刻画隐含波动率偏斜形态、捕捉非线性均值回归特征上的天然优势，其在加密货币市场的适配性研究显著增多。研究表明，加密资产在面临监管政策变动、宏观流动性冲击时，其波动率的均值回归表现出明显的非线性加速特征，这与 3/2 模型的数学性质高度契合（Fassas 等，2022）。此外，针对加密市场特有的 $7\times24$ 不间断交易机制，利用 3/2 模型结合高频数据进行实时波动率预测与期权动态对冲，已成为当前该领域的研究热点。但总体而言，现有关于 3/2 模型在加密市场的研究，多集中于静态参数校准与模型拟合效果对比，缺乏针对极端崩盘事件全周期的模型参数动态演化分析，也未对极端市场环境下不同定价方法的性能进行系统性实证检验，这一研究缺口为本文的实证分析提供了重要的研究空间。

\section{文献述评与本文研究定位}

综上所述，国际学术界围绕 3/2 随机波动率模型的理论构建、定价方法迭代与市场应用拓展已开展了丰富的研究，在模型解析性质、扩展框架构建、精确模拟方法优化、傅里叶定价方法改进等方面取得了一系列重要成果，为本文的研究奠定了坚实的理论基础。与此同时，国内研究多聚焦于 Heston 模型、SABR 模型等传统随机波动率模型，针对 3/2 模型的系统性研究相对匮乏，在极端市场环境下的定价方法性能评价、高波动加密资产市场的实证适配性检验等方面仍存在显著的研究缺口。

尽管现有研究已较为丰富，但仍存在三方面核心不足。第一，现有 3/2 模型定价方法的性能评价多局限于文献基准参数下的理想算例，缺乏覆盖常规市场与极端参数组合的全场景数值稳定性检验，部分方法在极端波动率的波动率、强相关系数、极短到期期限等真实市场常见参数下会出现数值溢出、计算失效等问题，方法的工程适用边界尚未得到清晰界定。第二，现有研究尚未基于真实极端市场校准参数完成定价方法的实证再检验，理论算例中的性能表现无法完全等价于真实市场中的实践效果，不同定价方法在极端行情下的精度、效率与鲁棒性差异仍缺乏实证依据。第三，3/2 模型在高波动加密货币市场的实证研究仍不完善，缺乏极端崩盘事件全周期下的模型参数动态演化分析与波动率微笑拟合效果检验，纯扩散 3/2 模型对极端市场的适配能力尚未得到充分验证。

基于现有研究的不足，本文明确了核心研究定位：以 3/2 随机波动率模型下的欧式期权定价为研究对象，系统构建三大范式十一类主流定价方法的统一实现框架，通过多场景数值实验与极端市场实证检验，全面评价各类方法的数值稳定性、定价精度与计算效率，厘清其工程适用边界；同时以比特币极端崩盘事件为样本，验证 3/2 模型在极端市场中的适配性，识别模型参数的时变规律。本文的研究旨在弥补现有研究在极端场景性能检验、真实市场实证分析等方面的不足，为 3/2 模型在高波动市场中的定价实践提供可靠的方法选择与实证支撑。